\chapter{知识演进的评价}
\label{chap:ket-evolution-evaluation}

前一章学会了“怎么把一次更新做出来”，但\textbf{“训练流程跑完、甚至新地区错误变少了”，\emph{远不等于}这个新版本可以放心上线}。因为一次更新可能同时在几个方向上表现不同，甚至互相打架。看三个真实的尴尬情形：
\begin{itemize}\item 经过几个季节的更新，\textbf{新地区错误减少了，可\emph{旧地区}成绩却\emph{跌破}了历史最好版本}（典型的“学会新的、忘了旧的”）；\item 负责人那条事实\emph{本身}改对了，但\textbf{依赖它的相关问答（谁签批、在哪个办公室）仍\emph{沿用旧答案}}（改了事实没改后果）；\item 撤回训练样本后，系统表面上不再提那个姓名了，但\textbf{\emph{“不再提”}并不能证明它\emph{参数里的训练影响}已经消除}。\end{itemize}
这些都说明：\textbf{没有任何单一数字能替你回答“更新成功没有}”。必须\textbf{\emph{分项}检查五件\emph{观察对象不同}的事}——当前任务表现、历史能力保持、后续学习能力、事实修改的正确性、以及训练样本撤回是否到位——本章就是建立这套“评价语言”，让你能有据可依地决定一个候选版本\textbf{接受、继续修、还是暂缓}。

本章在这一部分的位置是：第~\ref{chap:ket-reuse-selection}章解决“怎么\emph{比较复用收益}”，第~\ref{chap:ket-update-maintenance}章解决“怎么\emph{实施和单次验收}更新”；而系统是\emph{跨很多个阶段}持续演进的，本章就把“\textbf{\emph{跨阶段}比较所需要的指标、评价样本和\emph{持续协议}}”集中起来讲清楚。组织顺序对应一次严谨评价的思路：	extbf{第一，先给每一项主张规定“\emph{跟\emph{什么}参照}比”（11.1，没有参照的分数没有意义）；第二，再安排“样本的\emph{时间访问}和\emph{版本隔离}”，保证比较在因果上成立（11.2）；第三，最后学会“怎么\emph{联合解读}结果”，判断哪些证据足以支持接受、哪些要继续修、哪些该暂停（11.3）}。至于新任务迁移、多任务共享、来源融合的评价，由第三部分的第~\ref{chap:ket-evaluation}章统一展开，那里还有跨任务重复实验和统计区间的详细算法。

全章约定：风险、错误率\emph{越小越好}，准确率等成绩\emph{越大越好}，读公式时先看清方向。章内的花卉照片、阶段成绩矩阵、机构迁址/任命事实、费用数字都和前两章一样是\textbf{教学构造}，用途是让你\emph{亲手核对指标怎么算}，不能当成某方法的实测收益。还要再强调一次数据的三种身份（第 7、10 章反复出现）：\textbf{训练样本参与更新、验证样本用于选择配置（调强/早停）、独立测试样本只在模型\emph{锁定后}检查}；\textbf{同一张照片\emph{随着协议改变用途}时（比如延迟标签后来转去训练），必须\emph{重新声明}它还能支持哪一项评价}——身份不是一成不变的。

\section{评价目标与指标的确定}
\label{sec:ket-evolution-evaluation-metrics}

评价的第一件事不是“算分”，而是\textbf{\emph{在更新\emph{之前}}就写清楚“\emph{这次更新要满足哪些\emph{主张}}”}。接受一次更新，可能同时附带这些期望：当前错误要减少、旧任务的退步要受限、下一个新任务仍要能学会、指定的事实修改/撤回请求要真正落实。它们\textbf{可能同时成立，也可能\emph{相互冲突}}（新地区变准往往伴随旧地区风险）。所以要\textbf{事先列出两张清单：“需要\emph{改变}的行为”和“需要\emph{保持}的行为”}，并确立一条贯穿全章的铁律：\textbf{\textbf{这些主张\emph{不能用一个汇总平均分\emph{互相抵消}}——尤其“撤回合规、旧关键任务不退步”这类是\emph{硬要求}，不能因为另一项分数高就放行。}}本节先逐个建立度量这些主张的指标。

\subsection{当前任务表现与更新代价的度量}
\label{sec:ket-evolution-evaluation-current}

度量“当前任务有没有改善”，公平性设置与第 9 章收益判断完全一致：\textbf{先固定当前环境、物种含义、损失口径、评价样本的生成条件}，然后在\emph{同一批}测试样本上比较原版本和更新版本。方向沿用式~\eqref{eq:ket-reuse-risk-gain}（原风险 $-$ 更新风险，为正才算改善）；经验风险只代表“这批测试样本\emph{覆盖到的条件}”下的平均损失，不外推到没测的区域。要保留\textbf{两个不同的参照，缺一不可}：\textbf{原始版本}给更新前基线；\textbf{不更新基线}则让原权重走\emph{完全相同}的预处理、解码、调用接口。后者的作用是\textbf{防止“\emph{贪天之功}”}——如果你这次顺手修了图像解码的 bug 或换了预处理，变好可能来自这些\emph{外围修复}，而不是模型参数更新，同流程的不更新基线能把这部分贡献剥出来。

这里再次区分第 9 章那两种不能混用的主张，因为评价时它们对应不同数字：\textbf{若约束是“\emph{总共只给 80 个新标签}”，就比这 80 标签预算下\emph{能达到的成绩}；若要求是“达到某预定准确率”，就追踪“\emph{第一次在已测预算点达标}要花多少标签/算力”}。预算内始终没达标时，只能报告“已测到的最好成绩（观察上限）+ 未达标状态”，\textbf{\emph{不能}把你\emph{实际跑到底}花的成本冒充成“达标成本}”（没达标就没有达标成本）。流程纪律不变：按验证样本决定何时停，最终成绩仍要用\emph{独立测试}样本检查。

当前表现还要\textbf{\emph{分地区、分群组}报，不能只给一个总体平均}（11.3 会用例子证明平均方向会骗人）：新地区和仍要服务的旧地区分别计分，再按类别、遮挡情况、设备型号等列出分组表现。成本侧同样要把账列全（呼应第 9 章）：模型参数、回放样本、优化器状态、教师副本、恢复快照都占资源，且\textbf{标签数量、训练算力、持久存储、服务延迟、故障恢复等待时间要\emph{分项}报告}——\textbf{\emph{不能}用“新增了多少可训练参数”一个数字代表更新的总成本}（适配器参数少，但教师/回放/快照一样不少）。

第 10 章的 RanPAC 式“固定表示、解析求分类头”路线虽然不回放旧原始样本，但评价时别误以为它零成本、零风险：\textbf{它仍然要计算特征、累计 Gram 矩阵和各类统计、再求解分类头}，这些\textbf{统计量的字节数和矩阵求解的算力都\emph{不能填成零}}。验收要分两个不同问题：	extbf{第一，数值正确性}——增量累计的结果与“用同一固定岭目标、读全部已见样本\emph{联合}求解”的结果差多少（核对统计与求解有没有实现既定目标）；	extbf{第二，实际分类}——在“截至当前的\emph{共同类别集合}”上检查新类、旧类各自的错误（第 10 章二维例子显示统计无损时旧类仍可能因类别竞争变错）。前者一致\emph{不}代表后者好，\textbf{不能用“统计数值一致”代替“旧类能力保持}”。

\subsection{持续能力变化的度量}
\label{sec:ket-evaluation-continuous}

要判断“第三次更新有没有损害旧地区”，你需要的是\textbf{\emph{同一个任务在\emph{不同时刻}的成绩}，而不是一个最终总分}——最终总分高，可能是新任务的高分盖住了旧任务的下滑。于是持续学习评价有一个标准工具：\textbf{一张“阶段 × 任务”的\emph{成绩矩阵}}。设系统依次学 $K$ 个任务，阶段 $t$ 结束时模型是 $f^{(t)}$，它在任务 $k$ 那份\emph{固定}评价集 $Q_k$ 上的准确率记为 $a_k^{(t)}$。矩阵 $\bm A$ 里 $\bm A_{t,k}=a_k^{(t)}$，\textbf{行=评测时已经走到第几个阶段，列=在测哪个任务}；另外用\textbf{第 0 行记录“还没经历这串任务”的\emph{共同起点}}（学之前的基线）。读这张表的关键纪律是：\textbf{\textbf{只有\emph{沿同一列}上下比，才是在\emph{追踪同一项能力}随时间的变化}；沿一行横着比只是某时刻各任务的快照。}

先认识矩阵里三个区域（图~\ref{fig:ket-evaluation-matrix}会画出来）：\textbf{对角线 $a_k^{(k)}$}是“任务 $k$ \emph{刚学完那一刻}”的成绩，是后面很多指标的参照；对阶段 $t$ 而言，$k<t$ 的列是\textbf{历史任务}（学过、要检查忘没忘），$k>t$ 的列是\textbf{尚未学习的任务}（那一格本来还没法考试）。后者在某些协议下\textbf{仍可能有数字}——如果模型能借助已有输出接口对新任务\emph{零样本}预测（没训练也试着答），这正是前向迁移要看的；但要注意：\textbf{\emph{没测量}、或\emph{没有可用接口}时，那一格应保留“\emph{缺测}（--）”，\emph{绝不能填 0 分}}——“没考”和“考了个零蛋”是完全不同的信息，填 0 会严重拉低后续的平均值。

这套矩阵框架由第 10 章讲过的 GEM 类持续学习研究广泛采用：它从同一张表里\emph{分别}汇总三件事——最后一行的整体准确率、历史任务相对对角线的“后向变化”、以及学前相对基线的“学前变化”。本章统一用 $a$ 表示成绩（越大越好），以免和越小越好的风险 $R$ 混淆。一个容易被忽略的对照细节：如果你的两个待比方案\textbf{共享同一个预训练起点}，那么“学前对照”就应该取\textbf{这个起点\emph{尚未经历历史任务序列}时的成绩}，并明确说明\textbf{这\emph{不同于} GEM 原始论文的“\emph{随机初始化}基线}”——起点不同，比出来的东西也不同：这里比的是“在预训练之上，\emph{额外的历史经验}起了什么作用”\scite{ket-lopezpaz2017-gem}。

\subsubsection{历史能力保持与后向迁移的度量}
\label{sec:ket-evaluation-backward}

从矩阵能算出一系列指标，但\textbf{每个指标回答\emph{不同}问题、分母和参照\emph{不同}，不能互相替代}。第一个最直观：到阶段 $t$，模型已经承担了前 $t$ 项任务，它\emph{当下}在这些任务上的整体成绩如何？把\emph{当前这一行}里\emph{已经学过}的任务准确率等权平均，就是\term{阶段平均表现}（stage-wise average performance）：
\begin{equation}
  \bar a^{(t)}=\frac1t\sum_{k=1}^{t}a_k^{(t)}.
  \label{eq:ket-evaluation-stage-average}
\end{equation}
读这个平均有几个坑：分母 $t$ 数的是“\emph{已学习}的任务个数”，不是计划中的全部 $K$ 个、更不是照片张数。任务重要性不同时可以用预定权重，但要说明随任务陆续到达\emph{是否重新归一化}；想把未学任务也算进去的平均应另报。\textbf{最关键的局限}：\textbf{这条曲线的平均范围随 $t$ \emph{不断扩大}，所以一个\emph{新加入的困难任务}本身就会把 $\bar a^{(t)}$ 拉低}——平均下降\emph{可能}只是新任务难，\emph{未必}是旧能力忘了。\textbf{因此\emph{阶段平均分\emph{单独}不能}回答“旧能力到底降没降}”，那要靠下面固定任务列的指标。

要真正度量旧能力，必须固定任务列（盯住同一个任务 $k$），再选一个参照时刻。选不同的参照时刻，会得到两个名字相近、但含义和答案都可能相反的指标。第一是遗忘量（forgetting）$F_k^{(t)}$，它问当前成绩比这个任务历史上考过的最好成绩（峰值）低了多少，所以要先在“任务学完之后、到当前阶段之前”的所有检查点里取最大值当参照。第二是后向迁移（backward transfer，BWT）$B_k^{(t)}$，它问后续阶段的学习让这个任务比它刚学完那一刻（对角线）变好了还是变坏了，参照就是对角线成绩。直觉区别：F 的参照是历史最佳（你曾经考到的最高分），BWT 的参照是刚学完时（起点）。逐任务差值和对历史任务的等权平均定义为：
\begin{itemize}\item \term{遗忘量}（forgetting）$F_k^{(t)}$问：当前成绩比这个任务\textbf{\emph{历史上\emph{考过的最好成绩}（峰值）}}\emph{低了多少}？所以要先在“任务学完之后、到当前阶段之前”的所有检查点里取最大值当参照。\item \term{后向迁移}（backward transfer，BWT）$B_k^{(t)}$问：\textbf{\emph{后续}阶段的学习，让这个任务比它\emph{刚学完那一刻}（对角线）变好了还是变坏了}？参照就是对角线成绩。\end{itemize}
直觉区别：F 的参照是“\emph{历史最佳}”（你曾经考到过的最高分），BWT 的参照是“\emph{刚学完时}”（起点）。逐任务的差值和对历史任务的等权平均定义为：
\begin{equation}
  \begin{aligned}
    F_k^{(t)}&=\max_{k\leq u<t}a_k^{(u)}-a_k^{(t)},\\
    \bar F^{(t)}&=\frac1{t-1}\sum_{k=1}^{t-1}F_k^{(t)},\\
    B_k^{(t)}&=a_k^{(t)}-a_k^{(k)},\\
    \operatorname{BWT}^{(t)}&=\frac1{t-1}\sum_{k=1}^{t-1}B_k^{(t)}.
  \end{aligned}
  \label{eq:ket-evaluation-forgetting-bwt}
\end{equation}
读这两个量的符号和边界：$u$ 只遍历\emph{已保存}的历史检查点、不含当前阶段。$F>0$ 表示当前低于历史峰值（确实从最好成绩退了）；$B>0$ 表示当前比刚学完时还好（后续学习\emph{反过来帮到}了旧任务）。注意\textbf{当前若创出新高，$F$ 可以为\emph{负}}（现在比历史峰值还高），所以“负遗忘”不是笔误，它和正 BWT 一样表示改善。若你只关心“损失/退化”这一面，可以另报截断量 $\max(0,F)$（只惩罚真正退步），但要注明这是截断版、不能冒充原定义。两个边界情形：阶段 1 还没有历史任务，平均分母为 0，应记“\emph{未定义}”而不是 0；以及\textbf{历史检查点存得越密，越可能捕捉到一个更高的峰值}，所以跨方案比遗忘量时，\textbf{测量频率（多久存一次检查点）必须\emph{匹配}}，否则一个方案分高只是因为它存得勤。

光有定义还不够，下面用一个\textbf{可以亲手算}的三任务例子把这些指标全过一遍。把同一个“两物种识别”要求放进三个\emph{先后到达}的环境：$\mathcal T_1$=地区甲春季、$\mathcal T_2$=地区甲夏季、$\mathcal T_3$=地区乙秋季。先澄清：\textbf{这里的编号是“\emph{学习的先后顺序}”，和第~\ref{chap:ket-sharing}章那套“物种、病害、定位”三个\emph{不同}任务\emph{没有关系}}（恰好都用 $T_k$ 记号而已）。为保证可比，三个任务标签含义、输出空间都不变，模型也\emph{收不到任务身份}；每个任务\emph{另留 100 张固定照片}只做评测、不进训练。

图~\ref{fig:ket-evaluation-matrix}~给出人工正确数，每列分母均为100，因此格内数值也等于准确率百分数。共同起点在三项任务上都是50\%。阶段1没有评价任务3，该格保留缺测；右侧同列比较指出两个历史参照怎样产生不同答案。

\begin{figure*}[htbp]
  \centering
  % Teaching data, not measurements. Every task has 100 fixed evaluation images.
% Rows: (50,50,50), (80,55,missing), (86,82,60), (83,78,88).
\begingroup
\setlength{\tabcolsep}{2.5pt}
\setlength{\fboxrule}{0.6pt}
\renewcommand{\arraystretch}{1.65}
\begin{tikzpicture}[
  x=1mm,y=1mm,
  font=\figurefont\small,
  text=BookInk,
  cell/.style={minimum width=12mm,minimum height=7mm,inner sep=1mm,
    draw=BookMuted,line width=0.6pt},
  flow/.style={-{Latex[length=2mm]},draw=BookMuted,line width=0.8pt}
]
  \node[anchor=north west,inner sep=0pt,font=\figurefont\small\bfseries]
    at (0,4) {(a) 阶段 $\times$ 任务};
  \node[anchor=north west,inner sep=0pt] at (0,-3) {
    \begin{tabular}{@{}p{26mm}*{3}{>{\centering\arraybackslash}p{11mm}}@{}}
      \toprule
      阶段 & $\mathcal T_1$ & $\mathcal T_2$ & $\mathcal T_3$ \\
      \midrule
      0：共同起点 & 50 & 50 & 50 \\
      1：学完$\mathcal T_1$ &
        \fcolorbox{BookTeal}{BookTeal!10}{80} &
        \colorbox{BookAmber!12}{55} &
        \colorbox{BookAmber!12}{\textemdash} \\
      2：学完$\mathcal T_2$ &
        \colorbox{BookBlue!12}{86} &
        \fcolorbox{BookTeal}{BookTeal!10}{82} &
        \colorbox{BookAmber!12}{60} \\
      3：学完$\mathcal T_3$ &
        \colorbox{BookBlue!12}{83} &
        \colorbox{BookBlue!12}{78} &
        \fcolorbox{BookTeal}{BookTeal!10}{88} \\
      \bottomrule
    \end{tabular}
  };
  \node[anchor=west,inner sep=0pt,font=\figurefont\footnotesize]
    at (0,-49) {$k<t$：历史\quad $k=t$：当前\quad $k>t$：未学};
  \node[anchor=west,inner sep=0pt,font=\figurefont\footnotesize]
    at (0,-56) {单位：准确率（\%）；\textemdash：未测};

  \node[anchor=north west,inner sep=0pt,font=\figurefont\small\bfseries]
    at (81,4) {(b) 同列、不同时刻的比较};
  \node[anchor=west,inner sep=0pt] at (81,-6)
    {遗忘：任务1，历史最佳减当前};
  \node[cell,fill=BookBlue!12] (peak) at (88,-16) {86};
  \node[cell,fill=BookBlue!12] (nowf) at (120,-16) {83};
  \draw[flow] (peak.east) -- (nowf.west);
  \node[anchor=west,inner sep=0pt] at (131,-16) {$F_1^{(3)}=3$};

  \node[anchor=west,inner sep=0pt] at (81,-28)
    {BWT：任务1，当前减刚学完};
  \node[cell,draw=BookTeal,fill=BookTeal!10] (learned) at (88,-38) {80};
  \node[cell,fill=BookBlue!12] (nowb) at (120,-38) {83};
  \draw[flow] (learned.east) -- (nowb.west);
  \node[anchor=west,inner sep=0pt] at (131,-38) {$B_1^{(3)}=3$};

  \node[anchor=west,inner sep=0pt] at (81,-50)
    {零样本FWT：任务2，学前减起点};
  \node[cell,fill=BookPaper] (base) at (88,-60) {50};
  \node[cell,fill=BookAmber!12] (before) at (120,-60) {55};
  \draw[flow] (base.east) -- (before.west);
  \node[anchor=west,inner sep=0pt] at (131,-60) {$Z_2=5$};
\end{tikzpicture}
\endgroup

  \caption{三任务持续学习的教学矩阵。行是完成的阶段，列是评测任务；每格按100张固定照片计算准确率（\%）。实线框标出刚学完的对角线，浅底区分历史与未学区域，不编码成绩高低；\textemdash 表示未测。右侧保持任务列不变，分别以历史峰值、刚学完成绩和阶段0起点为参照计算遗忘、BWT与零样本FWT。}
  \label{fig:ket-evaluation-matrix}
\end{figure*}

先算阶段平均分（只平均\emph{已学}任务）：阶段 1 只有任务 1，$\bar a^{(1)}=80\%$；阶段 2 有任务 1、2，$\bar a^{(2)}=(86+82)/2=84\%$；阶段 3 三项都学了，$\bar a^{(3)}=(83+78+88)/3=83\%$。注意它从 84 微降到 83，但如前所述，这\emph{主要}是因为纳入了分数结构不同的新任务，单凭这个下降不能说旧能力退化。

到\textbf{阶段 2}，只有任务 1 是历史任务：它从刚学完的 80\% 升到了 86\%（学任务 2 反而帮到了任务 1）。代入定义：$F_1^{(2)}=80-86=-6$ 个百分点，$B_1^{(2)}=86-80=+6$ 个百分点。这里你能直观看到\textbf{“负遗忘”和“正后向迁移”是\emph{同一件“成绩创新高”的事实}、只是参照不同}；而且无论 $-6$ 还是 $+6$，都\emph{只描述了现象}，并没有\emph{解释}为什么会改善（可能是共享特征更稳，也可能是数据更丰富，要机制得另做对照）。

\begin{samepage}
阶段3要比较两个历史任务。任务1的历史峰值86\%降至83\%，因此$F_1^{(3)}=\max(80,86)-83=3$个百分点；任务2从刚学完的82\%降至78\%，因此$F_2^{(3)}=82-78=4$个百分点。平均遗忘为$(3+4)/2=3.5$个百分点，分母只包括这两个任务。若把参照改为各自刚学完的成绩，得到：
\[
  \operatorname{BWT}^{(3)}
  =\frac{(83-80)+(78-82)}2
  =-0.5\text{个百分点}.
\]

\end{samepage}
这个例子最有教学价值的地方在于：\textbf{任务 1 “\emph{比峰值低 3 个点}”（有所遗忘）和“\emph{比刚学完高 3 个点}”（后向获益）竟然\emph{同时成立}}！这说明 F 和 BWT \textbf{测的不是一回事，\emph{不能}用一个代替另一个、也\emph{不能}让它们互相抵消后说“没事}”。把两者的定义相加，中间的当前成绩 $a_k^{(t)}$ 恰好消掉：
\[
  F_k^{(t)}+B_k^{(t)}
  =\max_{k\leq u<t}a_k^{(u)}-a_k^{(k)}\geq0.
\]
右侧非负，是因为取最大值的候选里\emph{包含}了刚学完那一刻 $u=k$（峰值至少不低于起点）。这个恒等式也点明：只有当“刚学完的成绩\emph{正好就是}历史峰值”（中间没再创新高）时，才有 $F_k^{(t)}=-B_k^{(t)}$；在我们的例子里任务 1 曾在阶段 2 达到过 86\% 的更高峰，所以 F=3、BWT=3，\emph{并不}互为相反数。一句话：\textbf{报告持续学习结果时，F 和 BWT 应当\emph{各自}列出}。

逐列求差还有个隐含前提：\textbf{上下两个格子比的必须是“\emph{同一个\emph{仍然有效}的目标}}”，否则相减没有意义。最典型的是带时间的事实：还是第~\ref{ex:ket-overview-dated-fact}例里 2026 年 7 月 1 日甲所由甲交接给乙——评价“\emph{历史}任职”时本来就应答甲，评价“交接后的\emph{现任}”才应答乙。所以\textbf{更新后“现任答乙、历史仍答甲”是\emph{正确}的，把历史题从答甲“纠正”成答乙\emph{不是}改善，反过来把现任还答甲改成乙\textbf{也\emph{不是}“有害遗忘}”}。题目和标准答案都必须\emph{带时间条件}，历史问题、当前问题、不受影响的问题\textbf{分开计分}。同类地，若物种分类\emph{标准本身}变了（合并/拆分），应保留标准版本号或建立明确映射，\textbf{不能直接拿“语义已经不同”的两次成绩相减}。

同一批旧照片还可能走不同的接口被测，接口不同、任务难度就不同，报告时必须写明。第一，如果提供任务身份（告诉系统这是认种任务）让它选对应的旧任务头，评价的是一个带路由器的多头系统。第二，如果不给身份、让所有类别在统一输出空间里一起竞争，难度明显更高（模型还得自己判断这是哪类任务）。第三，如果允许调用历史模型检查点（旧任务直接用当时的旧模型答），那这个保持成绩属于“当前模型加版本库组成的整个系统”，已不再只属于当前这一个模型。所以报告要列清：给没给任务身份权限、实际路由选对的比率、能访问哪些检查点、以及版本库的存储费用，否则两个旧任务准确率不可比。

\subsubsection{前向迁移与持续可学性的度量}
\label{sec:ket-evaluation-forward}

矩阵还能回答一个朝\emph{前}看的问题：\textbf{历史经验有没有\emph{帮到\emph{后来}}的任务？}这里必须先把两件\textbf{发生在不同时刻}的事分开：\textbf{一是“\emph{学习之前}，模型\emph{零样本}就能把新任务预测得多好”（先验有没有用）；二是“\emph{学习过程中}，达到同等水平能\emph{省多少}代价”（学得快不快）}。旧任务保持得好，对这两个问题\emph{一个都没回答}。先看第一件：\term{零样本前向迁移}（zero-shot forward transfer）比较的是\textbf{新任务\emph{尚未接受任何训练}时}的成绩，问的是“先前的经验能否\emph{直接、不经过微调}就改善预测”。以共同起点成绩$b_k=a_k^{(0)}$为参照，将任务$k$紧邻学习前的成绩减去它，再平均任务2至$K$的差值：
\begin{equation}
  Z_k=a_k^{(k-1)}-b_k,\qquad
  \operatorname{FWT}_0
  =\frac1{K-1}\sum_{k=2}^{K}Z_k.
  \label{eq:ket-evaluation-zero-forward}
\end{equation}
读法和边界：任务 1 之前没有本序列的任何经验，所以只平均任务 2 到 $K$，分母 $K-1$（要求 $K>1$）。$Z_k>0$ 表示还没学任务 $k$，光靠已有经验就已经比共同起点强。这个比较要求两侧的标签含义、类别描述、预测接口完全相同；一个常见的“伪零样本”要排除：\textbf{如果你\emph{必须先用新任务标签拟合一个输出头}才能产出预测，那就\emph{不存在}该接口下的零样本成绩，应记“不适用（N/A）”，\emph{不能}把拟合之后得到的分数\emph{补进}零样本指标}（那已经用了新任务监督）。

图~\ref{fig:ket-evaluation-matrix}~中，任务2与任务3紧邻学习前分别为55\%和60\%，因此$Z_2=55-50=5$、$Z_3=60-50=10$个百分点，平均为$\operatorname{FWT}_0=(5+10)/2=7.5$个百分点。公式不需要阶段1的任务3成绩，故该处缺测不妨碍计算。若改问任务1的学习单独怎样影响任务3，就需要补测该格，不能由阶段2的60\%还原它。

读文献时要特别小心“同名不同义”的陷阱：CLEAR 基准补充材料里也有叫 Backward/Forward Transfer 的指标，但\textbf{它用的是\emph{绝对准确率}、且\emph{没有减去基线}}，和本章的“相对差值”口径不是一个量，照搬数字会误判。本例下三角为$(86+83+78)/3\approx82.33\%$，与本章$-0.5$个百分点的BWT不是同一个量；上三角缺少$a_3^{(1)}$，不能完整求平均。其Next-Domain Accuracy只取紧邻上对角线，本例为$(55+60)/2=57.5\%$，而非本章减去50\%起点后的7.5个百分点。使用文献名称时必须同时给出所取矩阵区域、运算和单位\scite{ket-lin2021-clear}

零样本成绩只回答“不学能不能猜个大概”，但更常见的实际问题是“学起来快不快”——达到同一水平，有历史经验的模型能不能少用标签或算力？这就要比较两条适应（学习）曲线。为公平，设 $c_{\star,k}^{\mathrm{hist}}$ 和 $c_{\star,k}^{0}$ 分别是：带历史经验的模型、和一个匹配好的无历史对照，在任务 $k$ 上达到同一个预定阈值 $\star$ 所需的适应成本（标签数等）。若两者实际都过线、且对照成本为正，用节省量除以对照成本，就得到“相对改善率”：
\begin{equation}
  E_k=\frac{c_{\star,k}^{0}-c_{\star,k}^{\mathrm{hist}}}
             {c_{\star,k}^{0}}.
  \label{eq:ket-evaluation-learning-cost-gain}
\end{equation}
$E_k>0$表示省成本，$E_k<0$表示额外花费；归一化参照是$c_{\star,k}^{0}$，不是两方案的成本之和。为与矩阵对应，构造预算点$n=(20,40,80,160)$：任务2的历史模型与无历史对照成绩分别为$(82,90,92,93)\%$和$(70,82,90,92)\%$，任务3分别为$(76,88,90,92)\%$和$(65,76,85,90)\%$。各预算点采用嵌套支持集、同一每类抽样规则及100张固定评测照片，并匹配容量与优化预算。90\%阈值下，任务2的到达量为40与80，任务3为80与160，所以$E_2=(80-40)/80=50\%$、$E_3=(160-80)/160=50\%$。

要理解这些扩展学习曲线是怎么来的、为什么不污染原系统：原序列里学任务 2 实际只用了 20 个标签、学任务 3 用了 40 个，所以它们在\emph{原序列}阶段终点的成绩就是图里的 82\% 和 88\%；而上面那些延伸到 160 标签的曲线，是\textbf{把阶段 1、阶段 2 的模型\emph{各复制一个独立副本}、在\emph{隔离的“探测”环境}里\emph{追加}预算跑出来的}。关键纪律：\textbf{\textbf{这些\emph{探测副本}和它们得到的反馈\emph{不返回}原序列}——否则你就等于借“评价”之名给原系统\emph{额外喂了训练数据}，历史成绩全被污染}。这样测到的是“\emph{同一个真实起点}\emph{之后}的学习效率”，是干净的；当然，探测本身烧的标签/算力\textbf{仍要登记为评价成本}（评价也不是免费的）。

又一次强调单位不可混：\textbf{学习成本节省 50\%（一个\emph{相对比例}）和零样本增益 7.5 个\emph{百分点}（一个\emph{绝对准确率差}）用的是不同单位，\emph{不能}把它们相加/相乘拼成一个所谓“前向迁移综合分}”。边界情形要会处理：任务 3 若到 160 标签仍没到 90\%，就照规则报告“$c_{\star,3}^{\mathrm{hist}}>160$”以及已测到达比例，\emph{不能编一个有限达标成本代进公式}；如果模型在零成本（零样本）时就已达标，直接报告零样本成功即可；而若\emph{对照}成本也是 0（两边都不用学），式~\eqref{eq:ket-evaluation-learning-cost-gain}分母为零、不能作除法。

还有一个更长期、也更容易被“旧分数还很好”掩盖的问题：\textbf{\emph{某一次}新任务学得快，\emph{不}代表系统在\emph{很多轮}更新之后\emph{仍然}学得动新东西}（第 8、10 章讲的可塑性，长期训练可能让模型僵化）。\term{持续可学性}（continual learnability）专门考察这种“\emph{后续学习的能力}”：做法是\textbf{从\emph{多个不同阶段}（第 1 阶段后、第 2 阶段后……）各取一个模型\emph{副本}，让它们去学\emph{同一个}隔离的探测新任务}，并\emph{严格固定}新样本、输出头、容量、调参预算和停止规则后比较学习曲线（探测反馈同样不交回原序列）。如果越后期的副本达到同样水平越慢，就说明可塑性在流失——哪怕它们的旧任务分数始终很好。

设三个历史阶段的旧任务成绩都为90\%，同一探测任务的过线步数却依次为40、80、160，匹配的无历史对照始终需要80步。相对成本改善便为$50\%,0,-100\%$。这个教学反例说明保持旧分数并不保证以后容易学习；它支持后续学习能力下降的行为描述，下降原因还需用新的对照检验。

做这种比较还有几个公平性陷阱：如果\textbf{后期模型\emph{扩容}了（参数变多），探测时必须\emph{补一个“同容量”的无历史对照}}，并列出新增存储/算力，否则它学得快可能只是因为模型更大，不能全记在“历史经验”头上；\textbf{任务\emph{到达顺序}的影响}要靠“换不同顺序重复跑”来检验，只换随机初始化覆盖不了顺序效应。要支撑“长期”主张，探测任务集还应\emph{包含旧环境再次出现、以及不同难度}的任务，并把历史训练与当前探测的开销\emph{分开记账}。最后记住证据边界：\textbf{哪怕一次适应很成功，结论也只\emph{限于那一次学习}，不能自动推广到无限未来}。

\subsection{事实修正与关联一致性的度量}
\label{sec:ket-evaluation-editing}

事实类修改的评价和分类任务很不一样，因为一条事实会被\emph{用在很多不同说法、不同推理里}。设想甲所迁到了另一座城市：直接问“它在哪”和问“出差该去哪个省”，考察的就是事实本身和它的\emph{连锁后果}。所以知识编辑评价要\textbf{把测试题\emph{刻意分成\emph{四组}、各管一件事}}：
\begin{itemize}\item \textbf{原请求题（直接题）}：用\emph{原请求的原句}问，检验“\emph{修改到底\emph{完成没有}}”；\item \textbf{等义改写/跨语言题}：用\emph{没参与编辑}的同义说法或别的语言问，检验“是不是只背下了那一个句子”（\emph{表达泛化}）；\item \textbf{依赖推理题}：答案需要\emph{用到新事实再往下推}（迁城后所属省份），检验“\emph{后果有没有\emph{传播}}”；\item \textbf{语义无关题}：问一些和这次修改\emph{毫不相干}的问题，检验“\emph{其他知识有没有被\emph{误伤}}（保持）”。\end{itemize}
四组要\textbf{各自报告成功率}，并注明是“按请求”还是“按问题”计分。如果各请求配的题数不一样，加权方式会严重影响结论：\textbf{若一条请求配了十种改写、另一条只有一种，直接按“题”求平均会让前者\emph{权重变成十倍}}。正确做法是\textbf{先在\emph{每个请求内部}求正确率，再对\emph{请求}\emph{等权}汇总}。

“答对没有”还依赖一套\emph{判分规则}，这套规则必须在\textbf{跑之前就\emph{冻结}}，不能看到结果再改：封闭答案可先按约定做 Unicode 规范化、统一大小写和空白，再匹配人工确认的等价答案集；日期、数值、单位用指定解析器和容差比较。\textbf{开放生成式回答}若用模型当裁判，要固定评分器版本、提示词、证据访问权限和标签体系，并在\emph{没参与方法选择}的人工子集上报告人-机评分一致性（否则自动评分不可信）。对“无法解析、模型拒答、多个答案互相冲突”也要预先规定怎么算分。两个常见的作弊/误判要杜绝：\textbf{只\emph{命中一个关键词}不能证明整句语义正确}；\textbf{事后看到某个有利答案才\emph{回写}判分标准}等于操纵评测。

为具体区分四组问题，给甲所增加一项独立于负责人交接的迁址请求：自2026年9月1日起从甲城迁至乙城，两城始终分别属于甲省与乙省。“9月以后所在地是哪座城市”及“9月应到哪座城市拜访”都应答乙城；“9月出差应前往哪个省”则需将新所在地接入“城市所属省”的关系，得到乙省。问2026年8月应去哪个省仍应答甲省；另一家未迁址机构的地址不应改变。时间与依赖关系由构造明确给出，不靠模型猜测。

构造10条编辑请求，每条配1道直接题、2道未用于编辑的等义改写和1道关联推理题。直接题答对9道、改写答对16道、推理答对6道时，请求内平均后再等权汇总，分别得到90\%、80\%与60\%；由于每条请求的题数相同，这里恰好也等于各组按题统计的比例。直接成功并未覆盖关联推理失败。

再设100道无关题，98道预测输出未变，预测输出保持率为98\%。它不等于正确率：原模型的错误也可能被保留。无关集应同时报告编辑前后准确率及预测变化率；若还检验概率分布变化，要预先定义距离和判断阈值。这样才能区分保护有效内容与仅维持原预测输出。

这里要精确理解“局部性”到底保护什么：迁址后所属省份跟着变，这是修改目标的合理后果，应该变、不算误伤；\term{局部性}（locality）要求保护的是“不依赖这条修改的内容”，而不是要求“所有字面上邻近的答案都不许变”——把该传播的也冻住反而是错的。MQuAKE 这类基准就是用多跳问题（沿事实链走好几步）来检验编辑后的关联后果，并区分反事实修改和真实时间更新两种设定。但要注意它的证据边界：它只覆盖前提正确、关系已知的有限长度事实链，测了这些链正确不能推广成“模型全部推理都一致”的保证\scite{ket-zhong2023-mquake}。

保持对象还包括不同层次的要求。物种分类成绩未降，不能说明历史负责人问答正确；事实问答正确，也不能说明未授权训练样本不会被模型复述。行为保留评价样本集应从修改目标与使用要求确定，不能只选容易保持的通用题。连续编辑后逐次复查仍有效请求、带时间条件的历史请求和相互依赖请求，才能追踪错误是否累积。

第 10 章的零空间编辑（AlphaEdit）在评价时要\textbf{把“\emph{内部保护}”和“\emph{端到端成功}”当\emph{两件事}分开测}：可以先在模型内部测“被保存旧键的线性响应变化有多大”，并报告投影阈值、近零方向的近似误差；但这只是内部指标，\textbf{还要在\emph{没参与}构造投影的事实题、改写题、依赖题上做\emph{端到端}预测检验}。\textbf{内部变化很小\emph{不能}证明端到端答案正确}（内部线性没动，非线性推理仍可能错）；而且如果连续编辑让\emph{上游表示}发生了漂移，当初用来保护的“旧键”可能已经不对应当前表示了，需要复查这些键是否仍有效。

\subsection{训练影响消除的度量}
\label{sec:ket-evolution-evaluation-unlearning}

如果请求不是“改所在地”，而是“\emph{撤回}某个曾用于训练的样本”，评价对象就发生了根本转变：\textbf{不再关心“某个答案内容对不对”，而关心“\emph{这条样本对模型的\emph{训练影响}还在不在}}”。这就是机器遗忘的评价，而它的\textbf{黄金参照非常明确}：\textbf{遗忘后的结果，应当和“\emph{假装这个样本\emph{从来没参加过}训练}”得到的结果\emph{对不上不上}}——即与“用\emph{剩余样本重新训练}”的模型对照。以$S$为原训练样本集、$D\subseteq S$为按样本标识指定的撤回集合，$\mathcal A$为原学习算法、$\mathcal U$为遗忘程序，两个模型写为：
\begin{equation}
  f_{\mathrm{unlearn}}=\mathcal U(\mathcal A(S),D),
  \qquad f_{\mathrm{retrain}}=\mathcal A(S\setminus D).
  \label{eq:ket-evaluation-unlearning-reference}
\end{equation}
由于训练本身带随机性（初始化、数据顺序等），严格的参照\textbf{比的不是“两次模型参数\emph{逐位相等}”（几乎不可能也没必要），而是“\emph{训练程序}和\emph{遗忘程序}\emph{各自产生的\emph{发布模型\emph{分布}}}\emph{不可区分}到什么程度}”。要区分两个层次：在\emph{有限输入}上比较输出概率分布，只能算一种\emph{行为层面的近似检查}，它\textbf{\emph{不等价于}模型分布层面的保证}（输出像，不代表内部机制/对其他输入的分布也像）。参照还要写明初始化、预处理、训练配置选择及受撤回训练样本影响的上游产物；从见过这些训练样本的检查点继续训练，不是当然合格的重训练。SISA通过限制训练样本的影响范围、重训受影响部分降低撤回开销，其实验同时报告准确率、重训时间与额外存储\scite{bourtoule2021sisa}

实践中可做的经验检查包括：在遗忘模型和重训练模型上比较剩余任务的效用、与撤回样本相关的行为、以及预定的隐私攻击结果是否接近。但有几个反直觉点必须记住：第一，撤回样本上的准确率不必然降到零，因为其他保留的合法样本可能独立支持同一个答案（第 10 章讲过，删的是这条样本的影响，不是这个事实）；第二，模型对有限提示拒答、或某一次成员推断攻击失败，都不能证明训练影响已消除（出口拦了或攻击没查到，不等于内部干净）。


\section{评价样本与持续协议的设计}
\label{sec:ket-evolution-evaluation-protocol}

指标定义得再严谨，\textbf{如果\emph{评价样本}本身不含你声称的那种变化，结论就是空中楼阁}。两个典型的“假评价”：\textbf{随机拆分\emph{同一季}的照片}，训练测试都来自同一季节分布，\emph{根本测不到}“真实季节\emph{先后到达}”这种时间漂移；\textbf{把旧照片\emph{反复混进训练}，又\emph{同时}把它们当\emph{独立}历史测试}，那旧任务考的其实是训练数据，保持率虚高。本节就解决“怎么\emph{设计样本和协议}，让评价真正\emph{覆盖}所声称的变化”，分三块：基准要覆盖什么序列/请求、怎么构造失败场景、以及样本用途/时间访问/版本怎么隔离。

\subsection{持续序列与修改请求的覆盖}
\label{sec:ket-evolution-evaluation-coverage}

要评价“长期使用”，数据集本身得有\textbf{时间戳和反馈先后次序}，否则无从构造序列。代表性的 CLEAR 基准利用 YFCC100M 图像的拍摄时间戳，构造了 2004 到 2014 年的视觉变化序列，主要在“类别集合固定、域随年份变”的设定下检验\emph{领域增量}。它用两种协议，对应两种问题：除了“在\emph{同一时间桶}内做训练/测试划分”，还有更贴近真实服务的\textbf{流式协议}——用当前模型先\emph{预测下一个时间桶}、\textbf{\emph{先封存}预测结果}，等标签到了再计分，然后这一桶才允许进入训练。流式协议检验的是“\emph{近未来}预测”能力；但要注意，一旦这些样本进入训练，它们就\textbf{\emph{不能再}充当独立的“历史保持”测试}（身份变了），所以若同时研究遗忘，得\emph{另留}一批历史审计样本并登记其额外成本。若同时研究遗忘，应另留历史审计样本集并登记额外资源。时间戳本身并未增加无限新类、无任务边界或事实撤回等条件。

连续预训练进一步涉及语料到达阶段与下游任务两种对象，应分别记录。每个语料阶段结束后，可以在固定下游任务上复测保持，再用隔离新任务测学习成本，同时计新增词元、旧语料回放及训练计算。新语料损失下降并不能说明下游能力保持；一次早期自监督预训练也不是长期持续预训练。CLEAR含有丰富无标签图像，但早期预训练的收益不能替代逐时间段更新后的这些检查。

对所在地这样的内容修改，评价样本应围绕事实依赖组织。MQuAKE-cf将事实链中部分内容改成反事实，MQuAKE-t考察真实时间变化；关系链其他部分保持不变，用多跳答案检验修改后果是否传播。它还包含按特定模型可回忆事实筛选的条件，需报告实体范围、链长、编辑数和问题生成方式。这组评价样本检验关联推理，而非仅检验编辑提示改写，也不是隐私删除或全部长期冲突的基准。

机器遗忘的训练与评价样本则必须对应实际训练样本、撤回粒度和时机，并有匹配的重训练参照，例如SISA在Purchase、SVHN等任务中的删除实验。随机训练样本撤回改为整类或整用户撤回，可能同时改变训练样本规模与任务要求，应另列这些条件，不沿用原来的删除结论。

这些基准各自控制不同条件。固定类别的时间流没有检验无限新类，已知关系的事实链没有覆盖全部推理，特定学习算法下的删除实验也不提供任意模型的删除保证。应分别报告阶段数、到达顺序、类别与环境变化、修改请求范围及撤回粒度，再将结果限定到这些条件。

\subsection{补充场景与失败样本的构造}
\label{sec:ket-evolution-evaluation-stress}

现成基准很少覆盖你真正担心的所有情形，所以要\textbf{\emph{按目标里“还没被覆盖的要求”\emph{主动补场景}}}。比如花卉系统要跨季节长期服务，而手头基准只测了“一次跨地区迁移”，那就还得补：稀有病害、严重遮挡、旧地区在新季节\emph{回归}、以及更长的任务序列；持续编辑场景则要补“同一事实\emph{再次}更新（乙又换成丙）、互相\emph{冲突}的请求、推理\emph{依赖的前提缺失}”这些压力情形。两条纪律：\textbf{补哪些场景由\emph{事先}的目标决定、\emph{纳入规则在看到结果前就\emph{固定}}（不能挑对自己有利的场景）；\emph{事后}才发现的失败可以整理成“诊断集”用于改进，但它\emph{不再能代表}“某个\emph{预先定义}总体的失败率}”（事后挑的样本有选择偏差）。

为了区分失败原因，可以构造成对条件。组合比较保持颜色、纹理、对象数量和背景分布，改变属性关系是否在训练中共同出现；分布适应固定标签与样本数，改变光照；关联编辑固定事实、模板和编辑数，改变依赖链长。若更长链也含更多冷门实体，就同时改变了熟悉度，不能把全部差异归给推理深度，应再匹配实体条件。

人为构造的对照场景\textbf{本身还得“\emph{有\emph{有效}语义}”，否则测出来的东西没有意义}：迁址请求若\emph{缺生效日期}，那历史问题和当前问题都\emph{没有唯一正确答案}，正确做法是补上日期、或把“明确报告证据不足”本身列为一个预期行为；冲突的两条请求必须给\emph{版本优先级规则}或“待核验”状态，不能让模型\emph{猜}该信哪条；合成的“病害”图案如果不符合真实成像和病理条件，那它只能用来测“模型对\emph{我定义的这个图案}作何反应”，\textbf{\emph{不能}解读为“临床/田间真实病害识别}”能力——别把玩具数据上的结论外推到真实任务。

\subsection{样本用途、时间访问与状态隔离}
\label{sec:ket-evolution-evaluation-access}

数据切分的“独立单位”必须是\textbf{\emph{对象/植株}，而不是文件}（呼应第 7、10 章）：同一植株连拍的一串照片、以及它们的裁剪/增强版本，要\textbf{归入\emph{同一个对象组}}，整组一起分到某个时期/环境，不能一组里几张进训练、几张进测试（否则答案几乎可互相推导）。三种用途各司其职：训练样本改参数、验证样本选配置、最终测试样本只在锁定后用；延迟标签后来进入训练，该样本对后续版本就不再是独立测试；用于修正事实、选超参的样本，同一轮里也不能再充当最终验收依据。

如果测试预测输入还会改变模型，隔离的不仅是样本，也包括状态。以$\sigma^{(t)}$表示模型参数、归一化统计、优化器、记忆、缓存及随机数等可变状态；独立任务或运行开始时恢复约定起点。同任务内跨样本是否保留、结束后是否丢弃，都需固定。只恢复权重却沿用上一任务统计或缓存，并未完成要求的重置。

在线/流式评价必须遵守一条因果铁律——\term{先预测后更新}（predict-then-update）：\textbf{用\emph{当前}状态对第 $i$ 个样本\emph{生成并\emph{封存}}预测 $\hat y_i$，\emph{等它的标签\emph{到达之后}}才能计分、并据此更新}，更新出的新状态只服务\emph{以后}的样本。若第 $i$ 个标签延迟了 $d_i$ 步才到，那么做 $\hat y_i$ 那一刻\textbf{只能读到“此前\emph{已经到达}”的反馈，\emph{绝不能}提前使用它自己的 $y_i$}（未来信息泄漏）。注意协议可以另行约定：TTA 里也有“先用当前无标签输入\emph{更新}再\emph{预测}”的玩法（先更新后预测），但必须明确写出无标签样本的范围、以及是否允许整批访问（整批看到的信息比逐条多）。

检测器同时参与维护时，保存当时的窗口、告警时点和随后执行的处置。ADWIN舍弃旧段是在改变所用近期统计，不能据此把窗口外错误从全程评价中删除。误报、漏报和检测延迟须按已知变化条件另报；若标签延迟，分别报告变化到可见反馈、可见反馈到告警的等待，不把反馈等待全归为检测器失效。逐时阈值的保证与整次运行的错误控制也要分开。

读到真实标签、更新后再对同一对象计分，得到的是训练拟合而非原来的预测成绩。预定在线规则允许流式标签训练未来状态，不允许事后替换已经封存的预测。

构造持续序列有三种来源，回答\emph{不同}问题、不能混用：	extbf{第一，真实时间序列}：必须严格保持因果访问（只能见过去），不能为了让平均成绩好看而任意打乱顺序；	extbf{第二，受控排列序列}：人为随机抽取多种任务/环境到达顺序、重复跑，专门检验“\emph{顺序效应}”（同一套任务换个顺序结果差多少）；	extbf{第三，旧环境回归序列}：让旧条件（比如温室季节）在后面再次出现，问系统“\emph{重新遇到旧条件时还能不能\emph{恢复}、并\emph{继续}学}”。三种结果应分别报告。如果还另设了独立的最终审计，审计样本\textbf{绝不进入}在线反馈流。

编辑划分则按请求及关系族隔离：原请求用于修改，改写、跨语言、关联题和无关题分别留出。重复直接答案的模板不构成多个独立事实成功；多跳测试可以提供被修改事实，不能同时把最终多跳答案作为示例。若要研究权限本身，可固定模型和样本，分别设置有无来源训练样本、任务身份或测试状态的组。它们测的是权限变化，强权限成绩不能称为同权限下的算法收益。


\subsection{比较基线、预算与重复单位}
\label{sec:ket-evolution-evaluation-controls}

评价一个“高级持续学习方法”，先要有足够强且预算公平的基线，否则收益站不住。最基本的两个对照是：朴素顺序训练（不做任何防遗忘处理，看裸奔会多糟），以及同等回放容量、同等算力预算下的简单回放（很多复杂方法其实赢不过调好的回放）。还可以放一个“能读取全部历史训练样本的联合训练”作强参照（它能看到所有旧数据，近似上限），但要明确：它的访问权限比你的方法更强，所以不能放进“同条件”排名里直接比；而且它只是一次有限优化得到的分数，也不能称为数学上的上界（没证明不存在更强解）。事实编辑比较原版本及相同请求下的直接微调；训练影响消除则比较指定学习算法排除撤回样本后的重新训练结果。

已有强预训练表示的分类协议，还可加入不继续改编码器的解析增量分类，检查更复杂保留机制是否值得新增计算。若由任务索引选择来源，应计入探针头拟合、梯度估计和索引准备，并用未参与索引准备及调参的任务检验选择后的真实效果，不能把嵌入距离当作收益。持续预训练比较则固定词元量、回放份额和可变模型参数，再分别改变学习率调度；同语料继续训练的对照有助于区分恢复调度与换域造成的变化。

检验保留约束时，固定骨干、训练样本与更新预算，比较有无该约束；移除回放时还要计入减少的样本访问和计算。把省下的预算用于普通更新，可以检查专门机制是否优于同量计算。候选配置、失败试训、教师查询、审计和恢复重算都计入费用，不能只留下最终成功版本的开销。

任何比较都要\textbf{\emph{重复}并报告波动}，但“\emph{重复什么}”决定了能排除哪种随机性，三者覆盖的东西不同：\textbf{重复随机\emph{初始化}}（同数据换种子重训）检验的是随机优化的波动；\textbf{重复\emph{任务顺序}}检验顺序效应；\textbf{\emph{重新抽取\emph{独立植株/任务}}}才真正\emph{增加对相应总体的覆盖}（样本层面的泛化）。要在\emph{可匹配条件}下做\emph{配对}比较（同一时期、同一任务、同一起点），报告配对差的均值和波动。一个经典错误要避免：\textbf{把\emph{同一张照片评价三次}当成“三张独立照片}”——它只是一次独立观测重复算了三遍，不会增加任何统计信息量。

\section{评价结果与维护决策的分析}
\label{sec:ket-evolution-evaluation-analysis}

\subsection{当前收益、历史保持与后续学习的联合判断}
\label{sec:ket-evolution-evaluation-tradeoffs}

11.3 要把各项结果合起来做维护决策，第一个陷阱就是：同一批结果，只因为“平均的单位”选得不同，连“改善了还是恶化了”的方向都可能反过来。看一个具体算例：地区甲有 900 张照片，准确率从 90\% 升到 92\%；地区乙只有 100 张，准确率却从 80\% 降到 70\%。按样本数加权平均（相当于按当前实际流量评价总正确数），从 $(810+80)/1000=89\%$ 变到 $(828+70)/1000=89.8\%$，上升 0.8 个百分点（甲样本多，主导了平均）。按地区等权平均（把两个地区当同等重要的服务对象），从 $(90+80)/2=85\%$ 降到 $(92+70)/2=81\%$，下降 4 个百分点。同一个更新，一个口径说变好、一个口径说变差：前者关心总共答对多少张，后者关心有没有哪个被平等承诺的地区被牺牲；再单看最差地区的 70\% 还能发现明显弱项。所以总体平均、地区等权、最差组、稀有类、关键任务的成绩都应并列报告，不能事后挑一个对自己有利的均值。

第 10 章讲过“保留”和“可塑性”可能此消彼长，评价时要把这种代价暴露出来：\textbf{一个保持约束可能既让旧任务更稳、又让后续适应变慢}。要干净地归因，就得做\textbf{控制变量的消融}：固定记忆、新训练样本、骨干网络，\emph{只改变约束强度}这一个因素，然后把“旧任务成绩、新任务学习曲线、总算力”三笔账\emph{并列}。如果同时又加了回放、又换了骨干，结果就\emph{无法}只归因于约束。要正确理解消融的证据强度：\textbf{消融能直接支持的是“\emph{你\emph{实际实施}的那套干预}对\emph{整个系统}的影响”；只有当\emph{伴随的其他差异都被控制住}，它才更接近“某个\emph{单一因素}的机制解释}”。

所以做联合判断时，几个量要\textbf{\emph{分别}保留、\emph{不许合并}成一个分}：当前风险（新环境准不准）、相对历史峰值的差（F）、相对刚学完的差（BWT）、后续学习成本（可塑性），各说各话。两个典型组合说明为什么不能二选一：\textbf{“历史成绩\emph{提高}、但后续学习\emph{变慢}”是\emph{两项\emph{同时发生}的不同变化}（BWT 正、可塑性下降），不能因前者高兴而忽略后者；“平均风险\emph{更低}、但某个\emph{必须服务}的地区\emph{退步}”也\emph{不能自动接受}（硬要求没过）}。图~\ref{fig:ket-update-evaluation}把“独立复测旧能力”和“隔离副本测后续学习”画成两个不交叉的分支，正是为了防止你用探测时的额外训练去污染、覆盖原版本的保持结果。

\subsection{事实修正、撤回要求与结论范围的核对}
\label{sec:ket-evolution-evaluation-decisions}

针对不同更新类型，决策时要回到各自的专门清单：事实修正逐项核对“原问法、等义改写、关联后果、无关行为”四组；撤回样本则另查“剩余任务表现、重训练参照、以及观察者可见的模型分布”。第 10 章反复强调的三个“\emph{不能替代}”在这里再确认一次：\textbf{外部事实条目变得不可检索、模型对相关问题拒答、有限次攻击未检出——这三者\emph{任何一个}都\emph{不能}替代“训练影响已消除”的证明}。还有一个回退陷阱：恢复旧模型时\textbf{必须\emph{继承}那些已生效的编辑与撤回（在旧快照上补做），\emph{不能}因为“旧任务分数较高”就\emph{跳过}这些合规检查}。

下结论时要诚实标注\textbf{证据的\emph{适用范围}和\emph{强度}}。所有持续保持/可学性的结论都\textbf{限于你\emph{实际}跑的序列长度、数据到达规则和容量}：三阶段的教学矩阵只是用来\emph{解释指标怎么算}，\emph{不能}证明系统能无限增长任务；固定类别的时间流没有检验开放世界里“无限新类”的情况；几条事实链答对了，也不代表所有隐含依赖都更新了。还要区分三种\emph{强度不同}的证据：\textbf{“未检出失败”（可能只是没测到）＜“在指定测试集上达标”（经验性的）＜“在明确假设下获得证明/保证”（最强但条件苛刻）}。写报告时\textbf{差异幅度、不确定性（区间）、没达标的运行、以及没测的范围，都要\emph{保留并如实写出}，不能只报最好的一次}。

最后给一条明确的\textbf{决策规则}：只有当候选\textbf{满足了\emph{事前}约定的“当前表现 + 保留要求”，并且事实修改、撤回、预算这些条件\emph{全部}通过}时，才能\emph{按其声明的范围}接受。不满足时有三条正当出路：证据不足就\textbf{补评价样本}、或\textbf{缩小结论的适用范围}、或\textbf{继续修正}；如果连一条\emph{合规的恢复路径}都没有（回退会带回已撤回内容），就应当\textbf{\emph{暂停}受影响的更新或服务}，而不是带病上线。一句话总结本章立场：\textbf{评价产出的是一份“\emph{有条件的维护决策}”，它\emph{并不要求}每次更新都硬凑出一个“成功”结论——“不通过/需补证据/暂缓”同样是正当而重要的评价结果}。

\section{本章小结}
\label{sec:ket-evolution-evaluation-summary}

全章可以收成一组“评价时必须分别回答的问题”：\textbf{当前任务改善了吗？旧能力是\emph{相对于\emph{哪个时刻}}保持住的（历史峰值还是刚学完）？后续任务还\emph{能不能继续学}？事实修改和样本撤回是否满足了它们\emph{各自的参照}（事实四组题、重训练/发布分布）？}而配套的那些指标——阶段平均分、遗忘量 F、后向迁移 BWT、零样本前向迁移 FWT、学习成本节省——\textbf{各用\emph{不同的分母和基线}，回答不同问题，\emph{绝不能}用一个综合分互相替代或抵消}（F=3 和 BWT=3 可以同时成立就是最直观的提醒）。

而这些数字\emph{可不可信}，取决于评价协议：样本的用途（训练/选择/验收）、时间访问（先预测后更新、标签延迟）、可变状态（归一化统计/缓存/优化器有没有越界共享）共同决定了比较是否成立。三条不可破的隔离是：\textbf{独立历史测试不进入在线反馈流、后续学习的隔离探测不回写原序列、重训练参照必须排除撤回样本及其派生物的依赖}。最后，把“逐任务的取舍、版本与事实的有效时间条件、全过程成本（含评价本身的开销）”放在一起核对，才能负责任地在“接受候选、继续修正、暂缓更新”之间做出决定。
